% =====================================================================
%  Phần C: Bài tập (P1: C1–C4, P2: C5)
% =====================================================================
\section{Bài tập và thực nghiệm}\label{sec:exercises}

% ---------------------------------------------------------------------
\subsection{C1. Giải tay SVM lề cứng}\label{sec:C1}
\owner{P1}
\textbf{Đề bài.} Lớp dương gồm $(2, 0)$ và $(0, 2)$, lớp âm gồm $(0, 0)$. Giải tay SVM lề cứng:
xác định $\vw$, $b$, độ rộng lề, các nhân tử Lagrange và kiểm tra toàn bộ điều kiện KKT.

\todo{C1: lời giải; đối chiếu với SMO của P2.}

% ---------------------------------------------------------------------
\subsection{C2. Bài toán XOR và ánh xạ đặc trưng}\label{sec:C2}
\owner{P1}
\textbf{Đề bài.} Chứng minh bài toán XOR không khả tách tuyến tính trong $\R^2$; tìm ánh xạ
đặc trưng tường minh $\phi$ với $\phi(\vx)\T\phi(\vz) = (\vx\T\vz + 1)^2$ và một siêu phẳng tách
trong không gian đặc trưng.

\todo{C2: lời giải.}

% ---------------------------------------------------------------------
\subsection{C3. Một hàm không phải hàm nhân xác định dương}\label{sec:C3}
\owner{P1}
\textbf{Đề bài.} Chứng minh $k(\vx, \vz) = \norm{\vx - \vz}$ không là hàm nhân xác định dương
bằng một phản ví dụ.

\todo{C3: lời giải.}

% ---------------------------------------------------------------------
\subsection{\texorpdfstring{C4. Ảnh hưởng của $C$ và $\gamma$}{C4. Ảnh hưởng của C và γ}}\label{sec:C4}
\owner{P1}
\todo{C4: thiết lập thí nghiệm; heatmap sai số kiểm định chéo lồng nhau;
số vector hỗ trợ và độ rộng lề theo $C$, $\gamma$; đối chiếu với cận ở Định lý~\ref{thm:B6-loo}.}

% ---------------------------------------------------------------------
\subsection{C5. Thời gian huấn luyện: SMO và hạ tọa độ đối ngẫu}\label{sec:C5}
\owner{P2}
\todo{C5: $n$ từ $10^3$ đến $10^5$, hàm nhân tuyến tính; đồ thị log-log; bảng thời gian.}
